{\begingroup
%%TITELSEITE%%
\iftablet
\ThisCenterWallPaper{1.2}{Strand.jpg} % XXX hier kannst du für die Tablet-Ansicht ein anderes Bild eintragen, nachdem du es in den Ordner img hochgeladen hast.
\else
\ThisLLCornerWallPaper{2.0}{Strand.jpg} % XXX hier kannst du für die Ansicht für andere Geräte als ein Tablet ein anderes Bild eintragen, nachdem du es in den Ordner img hochgeladen hast.
\fi
\noindent
\hspace{-8mm}\parbox{12cm}{\rule{10cm}{1.6pt}\vspace*{-\baselineskip}\vspace*{2pt}\\ % Dicke horizontale Linie
\rule{10cm}{0.4pt}\\[0.2\baselineskip] % Dünne horizontale Linie
\parbox[c][2cm]{10cm}{\centering\textsc{\LARGE \textbf{Tine De Booij}}\par}\\ % XXX hier änderst du deinen Titel!]

 \noindent\rule{10cm}{0.4pt}\vspace*{-\baselineskip}\vspace{3.2pt}\\ % Dünne horizontale Linie
\rule{10cm}{1.6pt}}\\[\baselineskip] % Dicke horizontale Linie
\endgroup}
\cleardoublepage
{\begingroup
% INNENTITEL
\centering 
\vfill
\parbox{\textwidth}{%
\rule{\textwidth}{1.6pt}\vspace*{-\baselineskip}\vspace*{2pt} % Dicke horizontale Linie
\rule{\textwidth}{0.4pt}\\[0.2\baselineskip] % Dünne horizontale Linie
\parbox{\textwidth}{%
\parbox[c][1.5cm][c]{1.5cm}{\includegraphics[height=1.5cm]{tikz/19}}
~ \hrulefill ~
\parbox[c][2cm][c]{7cm}{\centering\textsc{\Large \textbf{Het Leven van Tine De Booij}}}
 ~\hrulefill ~
\parbox[c][1.5cm][c]{1.5cm}{\scalebox{-1}[1]{\includegraphics[height=1.5cm]{tikz/19}}}}\\[0.2\baselineskip] % Title
\rule{\textwidth}{0.4pt}\vspace*{-\baselineskip}\vspace{3.2pt} % Dünne horizontale Linie
\rule{\textwidth}{1.6pt}}\\[\baselineskip] % Dicke horizontale Linie
\scshape % Small caps
Onderdeel van het familieverhaal \\[0.2\baselineskip] % XXX Kannst du löschen

\vspace*{2\baselineskip} % W
geschreven door \\[0.5\baselineskip]
{\Large Yvonne De Booij } \\
{\Large Anneflore Muller} \\
{\Large Dorus Faber} \\
{\Large Lisa van der Berg} \\
{\Large Jacqueline de Booij} \\


% XXX Hier kann ein zweiter Name eingetragen werden. Falls du keinen brauchst, kannst du diesen Abschnitt ab "\\" und bis hinter diesen blauen Text löschen.
\vfill
{Met mederwerking van  \\ {Tine de Booij}} % XXX Hier können weitere Namen eingetragen werden. Falls du keine eintragen willst, kannst du diesen Abschnitt ab "geschrieben von" und bis hinter diesen blauen Text löschen.
\vfill \vfill
{Met fotos van  \\ {Verscheidene Fotograven.}\par} %  XXX Hier kannst du deinen Fotografen eintragen oder den Abschnitt ab "{Fotos von" bis hinter diesen blauen Text löschen.

\vfill % 

{\scshape 2021} \\[0.3\baselineskip] % XXX Hier kannst du das Jahr der Veröffentlichung / Fertigstellung eintragen, oder die Zeile ab "{\scshape" bis hinter diesen blauen Text löschen. 

\endgroup}